\subsection{结构的弱化}

在本节中，字母 r、s、r'、s' 或表示整数 \(\geqslant1\)，或表示符号 \(\infty\) 和 \(\omega\) 之一。假设 K = R。

5.13.1. 设 \(r \leqslant s\)，且 X 是 C＾s 类流形。在拓扑空间 X 上存在唯一的 C＾r 类流形结构，使得原结构下 X 的每个图也是这一新结构下 X 的图。令 \(C^s\)\(C^r\)\(X_r\) 为如此定义的 C＾r 类流形。称它是通过弱化 X 的流形结构从 X 得到的，或称其流形结构蕴含于 X 的流形结构中。弱化的概念具有传递性：若 \(C^r\)\(r' \leqslant r\)，则有 \(X_{r'} = (X_r)_{r'}\)；它与乘积相容：若 Y 是 C＾s 类的，则有 \(C^s\)\((X \times Y)_r = X_r \times Y_r\)；在第 5.11.2 号的假设下，对于 \(X \times_s Y\) 也有类似结果。

令 \(a\in X\)，\(c\) 是 X 在 \(X\)\(a\) 处的一个图；\(c\) 也是 \(X_r\) 在 \(a\) 处的一个图。由此根据（5.5.1）得到同构

\[
\theta_ {c}: \mathrm{E} \rightarrow \mathrm{T} _ {a} (\mathrm{X}), \quad \theta_ {c} ^ {\prime}: \mathrm{E} \rightarrow \mathrm{T} _ {a} (\mathrm{X} _ {r}),
\]

从而得到同构 \(\theta_{c}^{\prime}\circ\theta_{c}^{-1}:T_{a}(X)\to T_{a}(X_{r})\)。该同构与 \(c\) 的选择无关；它使我们能够把 X 与 \(X_{r}\) 在 \(a\) 处的切空间等同起来。

5.13.2. 设 X（分别为 \(\mathbf{X}'\)）是 \(\mathbf{C}^s\) 类（分别为 \(\mathbf{C}^{s'}\) 类）流形，且 \(r\) 满足 \(r \leqslant \inf(s, s')\)。若映射 \(f: \mathbf{X} \to \mathbf{X}'\) 是从 \(\mathbf{C}^r\)\(\mathbf{X}_r\) 到 \(\mathbf{X}_r'\) 的态射，则称其为  类映射；这种映射对于所有  都是 \(\mathbf{C}^{r'}\) 类的。此外，\(r' \leqslant r\)\(f: \mathbf{X}_r \to \mathbf{X}_r'\) 在点 \(a \in X\) 处的切线性映射，与把 f 视为从 \(f\)\(\mathbf{X}_{r'}\) 到 \(\mathbf{X}_{r'}'\) 的态射时的切线性映射相同。

通常用同一符号表示流形 X 和 \(X_{r}\)；因此，若 X 是 \(C^{s}\) 类的，则“X 的 \(C^{r}\) 类子流形”（\(r \leqslant s\)）这一表述意指“\(X_{r}\) 的子流形”。

